\section{Background: The Biology of Extreme Radiation Resistance}
\label{sec:background}

\subsection{Ionizing radiation as a selective force}
Ionizing radiation damages cells through two routes. Direct deposition of
energy in DNA produces double-strand breaks (DSBs), the most lethal common
lesion; a single unrepaired DSB can be sufficient to kill a dividing cell.
Indirect damage is mediated by water radiolysis, which generates hydroxyl
radicals, superoxide, and hydrogen peroxide, and these reactive oxygen
species (ROS) oxidize proteins, lipids, and nucleic acids far from the site
of the original photon absorption. For desiccation-tolerant organisms the
indirect route dominates: at low water activity the primary lesion burden
shifts toward oxidative protein carbonylation, and survival correlates more
strongly with proteome protection than with DNA repair capacity alone.

\subsection{The canonical model: \emph{Deinococcus radiodurans}}
\emph{D. radiodurans} R1 survives acute gamma doses above 10~kGy without
loss of viability, roughly three orders of magnitude beyond the lethal dose
for \emph{Escherichia coli}. Three decades of work have established that its
resistance is not attributable to unusual target hardening of DNA. Instead,
a layered defense acts on the proteome and the metabolome: (i)~manganese
accumulation, with intracellular Mn(II) concentrations reaching millimolar
levels and Mn/Fe ratios far above those of sensitive bacteria; (ii)~Mn(II)
complexes with orthophosphate and small metabolites (peptides, nucleosides)
that scavenge superoxide and hydroxyl radicals non-enzymatically; (iii)~an
unusually deep enzymatic antioxidant battery, including multiple catalases,
superoxide dismutases, and peroxiredoxins of the AhpC/TSA family; and
(iv)~an efficient homologous-recombination machinery operating on a
multipartite, polyploid genome that reassembles shattered chromosomes from
overlapping fragments through extended synthesis-dependent strand
annealing. The resulting picture is that of a cell whose DNA is damaged at
the same rate as any other cell's but whose proteins remain functional
throughout irradiation, permitting rapid repair once conditions allow.

\subsection{Convergent resistance in other clades}
Extreme resistance has evolved independently in several lineages, and each
case is an independent natural experiment. Within bacteria, the
Deinococcaceae contain multiple resistant species (\emph{D. radiodurans},
\emph{D. geothermalis}, \emph{D. deserti}), and the family is treated in
this study as a single origin because of shared ancestry. Among animals,
tardigrades such as \emph{Ramazzottius varieornatus} tolerate kilogray
doses in the desiccated state, and at least one lineage-specific protein,
Dsup, binds nucleosomes and reduces DNA damage in heterologous systems,
including human cells. Bdelloid rotifers such as \emph{Adineta vaga}
combine desiccation tolerance with radiation resistance and carry heavily
reorganized genomes with signatures of horizontal gene transfer, including
genes of bacterial and fungal origin implicated in stress response. These
three origins - Deinococcaceae, Tardigrada, and Bdelloidea - are separated
by more than a billion years of evolution, so any molecular feature shared
across them is a candidate convergence on a common solution rather than a
shared inheritance.

\subsection{Why oxidative defense is the leading axis}
Across resistant taxa the recurring observation is elevated antioxidant
capacity: peroxiredoxins, glutaredoxin and thioredoxin systems, flavin
reductases that recycle redox cofactors, and metal-homeostasis systems
that limit iron-catalyzed Fenton chemistry. Peroxiredoxins of the
AhpC/TSA family are peroxide-scavenging enzymes with exceptionally high
catalytic rates; their up-regulation is a common response to irradiation
in bacteria, and their abundance correlates with desiccation tolerance in
tardigrades and rotifers. If radiation resistance across kingdoms shares a
single dominant mechanistic axis, comparative genomics should recover
orthogroups in this pathway as convergently retained or expanded in
resistant lineages while absent from sensitive relatives - the hypothesis
tested here.

\subsection{The specificity problem}
A central confound is that desiccation, oxidation, heat, and radiation
stresses are correlated in nature. Organisms that survive drying
necessarily solve the oxidative-damage problem, and radiation resistance
may in part be a byproduct of desiccation adaptation. Any genomic
enrichment found in resistant species could therefore reflect general
stress adaptation rather than radiation-specific selection. This study
addresses the confound directly with a pre-registered, stress-matched
control panel (Section~\ref{sec:stresspanel}): stress-tolerant but
radiation-sensitive species are subjected to the identical pipeline, and
candidate orthogroups must be absent from that panel to retain a
radiation-specific interpretation.
